\documentclass[11pt,a4paper]{article}
% Estilo comum dos PDFs individuais da Sprint 1.
\usepackage[utf8]{inputenc}
\usepackage[T1]{fontenc}
\usepackage[brazilian]{babel}
\usepackage[a4paper,margin=2.35cm]{geometry}
\usepackage{xcolor}
\usepackage{hyperref}
\usepackage{enumitem}
\usepackage{listings}
\usepackage{parskip}
\usepackage{needspace}
\definecolor{azul}{RGB}{20,63,112}
\hypersetup{colorlinks=true,linkcolor=azul,urlcolor=azul}
\setlist[itemize]{leftmargin=*,itemsep=3pt}
\lstset{basicstyle=\ttfamily\small,frame=single,breaklines=true,columns=fullflexible}
\sloppy
\newcommand{\codigo}[1]{\texttt{#1}}
\newcommand{\titulo}[2]{\begin{center}{\LARGE\bfseries #1\par}\vspace{4pt}{\large #2\par}\vspace{8pt}\rule{0.8\textwidth}{0.4pt}\par\vspace{5pt}{\small Matheus Moreli Durão Silva -- Matrícula 12621CBS217\par Cibersegurança - Bacharelado - Noturno\par}\end{center}}
\newcommand{\cardfonte}[1]{\par\smallskip\noindent\textit{Referência: }\href{#1}{Trello - enunciado e checklist do card}.}

\begin{document}\titulo{S1\_07 - Estruturas condicionais}{Classificar a severidade de uma vulnerabilidade}
\section*{Entregável correto}
O programa recebe a severidade de uma vulnerabilidade e mostra a classificação usando condicionais. As opções estudadas são Baixa, Média, Alta e Crítica; qualquer outro valor precisa ser identificado como inválido.
\section*{Como a decisão funciona}
\codigo{if} executa um bloco quando a condição é verdadeira. \codigo{elif} testa outra possibilidade apenas se as anteriores falharem, e \codigo{else} cobre o caso restante. \codigo{==} compara valores; \codigo{and} exige duas condições verdadeiras; \codigo{or} aceita ao menos uma. A indentação (espaços à esquerda) pertence à sintaxe do Python e define o bloco.
\begin{lstlisting}
severidade = input("Severidade: ").strip().lower()
if severidade == "baixa":
    print("Classificacao: Baixa")
elif severidade == "media":
    print("Classificacao: Media")
elif severidade == "alta":
    print("Classificacao: Alta")
elif severidade == "critica":
    print("Classificacao: Critica")
else:
    print("Severidade invalida.")
\end{lstlisting}
\codigo{strip()} remove espaços das pontas e \codigo{lower()} deixa a comparação menos sensível a maiúsculas. Isso é uma validação simples, não uma regra ``secreta''.
\section*{Aplicação no projeto}
O \codigo{app.py} usa decisões para validar campos, verificar se existe ativo e interpretar escolhas do menu. A classificação de severidade volta a aparecer nas vulnerabilidades; mais adiante, \codigo{Enum} reduz o risco de escrever valores diferentes para a mesma opção.
\section*{Checklist demonstrado}
Tomada de decisão, \codigo{if}, \codigo{if/else}, \codigo{if/elif/else}, operadores relacionais/lógicos e validação foram aplicados. Testar uma opção válida e uma inválida comprova os dois caminhos.
\cardfonte{https://trello.com/c/QvZa8m7a}
\end{document}
